% ======================================================================
% capitulos/linea3/capitulo12.tex
% ----------------------------------------------------------------------
% LÍNEA CE03 — INFRAESTRUCTURA TECNOLÓGICA
% Semestre 1 · Entregable 2 de 2: ANÁLISIS Y DISEÑO DE CENTRO DE DATOS HÍBRIDO
% Estructura según la plantilla oficial
% S1_E2_Analisis_y_Diseno_de_Centro_de_Datos_Hibrido.docx
% ======================================================================
\chapter{Análisis y Diseño de Centro de Datos Híbrido}
\label{chap:ce03-dc}

\entregable{CE03 --- Infraestructura Tecnológica · competencia CE033 Implementación de Centro de Datos}{Semestre 1 · Entregable 2: Análisis y Diseño de Centro de Datos Híbrido}{CE0331 Diseño de centro de datos}{Completo (diseño, sin implementación).}


Este capítulo es el segundo entregable del semestre~1 de la línea CE03. Diseña dónde
y cómo se alojan los servicios de WISP: un centro de datos local virtualizado con
\textbf{VMware ESXi} en la oficina central y un entorno en \textbf{AWS} para el
concentrador VPN, la contingencia y los respaldos. Se apoya en el diseño de red del
\cref{chap:ce03-red}, con las mismas redes y extremos.

% ======================================================================
\section{Resumen ejecutivo del diseño del centro de datos híbrido}

\begin{itemize}
  \item \textbf{Proyecto:} centro de datos híbrido del Sistema de Gestión Integral de
  WISP (Juliaca, Puno).
  \item \textbf{Problema:} hoy todo el sistema corre en una sola PC de oficina (AMD
  A8-7600, 4~GB de RAM, disco único) sin redundancia, sin virtualización y con la
  base de datos, la publicación y la mensajería en el mismo sistema operativo. La VPN
  depende de una VPS de terceros y los respaldos quedan en el mismo equipo.
  \item \textbf{On-premise (VMware ESXi):} un servidor de rack (8 núcleos, 64~GB, SSD en
  RAID~1) con ESXi~8, siete máquinas virtuales separadas por función y zona de
  seguridad. Se elige ESXi por ser el hipervisor que exige la rúbrica, su madurez en
  aislamiento de cargas y su disponibilidad en edición gratuita o académica.
  \item \textbf{Cloud (AWS):} una VPC en São Paulo con el hub WireGuard (necesita la IP
  pública fija que la oficina no tiene), una réplica de contingencia y respaldos en S3
  fuera de la oficina.
  \item \textbf{Dimensionamiento:} 19 vCPU, 35~GB de RAM y unos 630~GB asignados,
  frente a 16 hilos, 64~GB y 1,9~TB del host: un 39\,\% de margen en RAM y un 65\,\% en
  disco. En AWS, dos instancias \texttt{t4g.medium}.
  \item \textbf{Siguiente paso (semestre~2):} instalar el host, migrar los contenedores
  a las VM, desplegar la VPC y ejecutar al menos una prueba completa de recuperación.
\end{itemize}

\begin{figure}[htbp]
\centering
\begin{tikzpicture}[x=1cm,y=1cm]
  \node[caja,text width=4.6cm,minimum height=24mm] (onp) at (-3.6,0) {\textbf{On-premise (oficina central)}\\[2pt]Rack 18U · FW01 · SW-CORE\\ESXI-01: 7 VM (app, BD, DMZ,\\monitoreo, SIEM, IDS, señuelo)\\NAS-01 respaldos · UPS};
  \node[caja,text width=4.2cm,minimum height=24mm] (aws) at (3.8,0) {\textbf{AWS sa-east-1}\\[2pt]VPC 10.50.0.0/16\\EC2-WGHUB (hub VPN)\\EC2-DR (réplica)\\S3 (respaldos)};
  \draw[tunel] (onp) -- node[font=\tiny,above]{WireGuard} node[font=\tiny,below]{réplica · respaldos} (aws);
  \node[cajaext] (z) at (3.8,-2.6) {13 zonas (WireGuard)};
  \draw[tunel] (z) -- (aws);
\end{tikzpicture}
\caption{Topología ejecutiva simplificada del centro de datos híbrido.}
\label{fig:ce03-dc-ejecutiva}
\end{figure}

% ======================================================================
\section{Descripción de la organización y necesidades de infraestructura}

WISP es un ISP de fibra óptica con unos 455 clientes en 13 zonas. Su operación diaria
depende de un sistema de gestión propio (línea CE02) que registra clientes, pagos,
planta externa e inventario, y que gestiona los routers de las zonas.

\begin{table}[htbp]
\centering
\caption{Servicios que requieren cómputo o almacenamiento y entorno asignado.}
\label{tab:ce03-servicios-dc}
\small
\begin{tabularx}{\textwidth}{L{3.6cm}L{2.2cm}XL{1.6cm}}
\toprule
\textbf{Servicio} & \textbf{Entorno} & \textbf{Criterio de reparto} & \textbf{Info.} \\
\midrule
API, panel web, sitio, WhatsApp & On-premise & Uso diario desde la oficina; costo fijo; ya opera en local. & Interna \\
Base de datos PostgreSQL & On-premise & Contiene datos personales: queda bajo control directo de la empresa. & Confidencial \\
Publicación (Nginx, cloudflared) & On-premise (DMZ) & Publica sin IP pública mediante un túnel de salida. & Pública \\
Monitoreo, SIEM, IDS, señuelo & On-premise & Necesitan ver el tráfico de la oficina (espejo del troncal). & Interna \\
Concentrador VPN & AWS & Requiere IP pública fija y alta disponibilidad, que la oficina no tiene. & Interna \\
Réplica de contingencia & AWS & Debe estar fuera de la oficina para sobrevivir a un incidente local. & Confidencial \\
Respaldos fuera de sitio & AWS (S3) & Tercera copia de la regla 3-2-1, durable y económica. & Confidencial \\
\bottomrule
\end{tabularx}
\end{table}

\begin{table}[htbp]
\centering
\caption{Situación actual (AS-IS) y crecimiento estimado.}
\label{tab:ce03-asis-dc}
\small
\begin{tabularx}{\textwidth}{L{3.4cm}X}
\toprule
Servidor actual & PC AMD A8-7600 (4 núcleos), 4~GB de RAM, disco de 450~GB, Ubuntu Server 24.04, Docker Compose con 6 contenedores; consume 1,3 a 1,8~GB de RAM en reposo. \\
Publicación & Cloudflare Tunnel (sin puertos abiertos). \\
VPN & VPS de terceros con el hub WireGuard; algunas zonas aún usan SSTP. \\
Respaldos & \texttt{pg\_dump} diario en el mismo equipo, retención de 30 días, copia manual semanal a una laptop. \\
Crecimiento a 3 años & Del orden de 910 clientes y hasta 25 zonas; evidencias fotográficas de unos 30~MB por cliente (unos 27~GB). \\
\bottomrule
\end{tabularx}
\end{table}

% ======================================================================
\section{Alcance del diseño de centro de datos híbrido}

\subsection{Declaración de alcance}
El diseño comprende la infraestructura virtualizada VMware ESXi del centro de datos
local de la oficina central (rack, host, almacenamiento, red virtual y máquinas
virtuales) y los recursos en AWS (VPC, subredes, instancias EC2, Security Groups, NACL,
IAM y S3), incluidos su dimensionamiento, su integración y los mecanismos de
disponibilidad y respaldo.

\subsection{Límites}
\begin{itemize}
  \item \textbf{Físicos:} rack de 18U en la oficina central con FW01, SW-CORE, ESXI-01,
  NAS-01, UPS y PDU.
  \item \textbf{Tecnológicos:} ESXi~8, vSwitches, \emph{port groups}, VM; AWS VPC, EC2,
  EBS, S3, Security Groups, NACL, IAM y CloudWatch.
  \item \textbf{Organizacionales:} el \emph{responsable de infraestructura} administra
  el rack, ESXi y AWS; el \emph{responsable de sistemas}, las VM APP-01 y DB-01 y los
  respaldos de datos; la gerencia aprueba los accesos a la cuenta de AWS.
\end{itemize}

\subsection{Exclusiones}
\begin{itemize}
  \item \textbf{Equipamiento de las zonas (routers, OLT, gabinetes):} se trata en el
  diseño de red; no aloja servicios del centro de datos.
  \item \textbf{Estaciones de trabajo del personal:} no alojan servicios.
  \item \textbf{Alta disponibilidad con dos hosts y vCenter:} se descarta por costo; se
  compensa con la réplica en AWS.
  \item \textbf{Cloudflare:} servicio gestionado por el proveedor; solo se diseña su uso.
\end{itemize}

% ======================================================================
\section{Objetivos del diseño de centro de datos}

\begin{table}[H]
\centering
\caption{Objetivos medibles del diseño de centro de datos.}
\label{tab:ce03-objetivos-dc}
\small
\begin{tabularx}{\textwidth}{XL{3.6cm}L{2.2cm}L{2.2cm}}
\toprule
\textbf{Objetivo} & \textbf{Indicador} & \textbf{Meta} & \textbf{Responsable} \\
\midrule
Dimensionar el host ESXi & Recursos asignados frente a requeridos & 100\,\% cubierto con $\geq 30$\,\% de margen & Infraestructura \\
Separar cargas por función y zona de seguridad & Servicios en VM dedicadas & 7 VM, 100\,\% & Infraestructura \\
Definir una arquitectura AWS costo-eficiente & Servicios seleccionados y justificados & 100\,\% documentado & Cloud \\
Planificar la disponibilidad mínima & RTO/RPO por servicio definidos & RTO $\leq 2$~h en servicios críticos & Infraestructura \\
Respaldos fuera del sitio & Copias en S3 & 3-2-1 para BD y evidencias & Sistemas \\
\bottomrule
\end{tabularx}
\end{table}

% ======================================================================
\section{Arquitectura on-premise --- diseño (VMware ESXi)}

\subsection{Topología física propuesta}

\begin{figure}[H]
\centering
\begin{tikzpicture}[x=1cm,y=1cm]
  \tikzset{u/.style={draw=upeublue,fill=upeublue!6,text width=5.45cm,align=center,font=\scriptsize,anchor=north west,inner sep=1.5pt}}
  \draw[upeublue,very thick] (-0.15,0.15) rectangle (5.75,-9.15);
  \node[u,minimum height=5mm] at (0,0) {U18 · Patch panel 24 puertos (Cat6)};
  \node[u,minimum height=5mm] at (0,-0.5) {U17 · Organizador de cables};
  \node[u,minimum height=5mm,fill=red!6,draw=red!60!black] at (0,-1.0) {U16 · FW01 RB5009 (bandeja)};
  \node[u,minimum height=5mm] at (0,-1.5) {U15 · SW-CORE CRS326};
  \node[u,minimum height=5mm] at (0,-2.0) {U14 · Organizador de cables};
  \node[u,minimum height=10mm] at (0,-2.5) {U12--13 · Libre (segundo host)};
  \node[u,minimum height=5mm,fill=upeublue!18] at (0,-3.5) {U11 · ESXI-01 Dell PowerEdge R250};
  \node[u,minimum height=10mm] at (0,-4.0) {U9--10 · NAS-01 DS224+ (bandeja)};
  \node[u,minimum height=15mm] at (0,-5.0) {U6--8 · Ventilación};
  \node[u,minimum height=5mm] at (0,-6.5) {U5 · PDU con protección};
  \node[u,minimum height=10mm,fill=upeugray!15] at (0,-7.0) {U3--4 · UPS APC 1500 VA (2U)};
  \node[u,minimum height=10mm] at (0,-8.0) {U1--2 · Reserva (baterías)};
  \node[cajaext,text width=4.4cm,align=left] at (9.0,-2.0) {\textbf{Enlaces}\\FW01 SFP+ $\leftrightarrow$ SW-CORE (troncal)\\ESXI-01: 2 NIC gestión + 2 NIC datos\\NAS-01: 1 NIC en VLAN 20\\Puerto espejo $\rightarrow$ ESXI-01 (IDS)};
  \node[cajaext,text width=4.4cm,align=left] at (9.0,-6.4) {\textbf{Energía}\\UPS $\rightarrow$ PDU $\rightarrow$ todos los equipos\\BIOS: encendido tras corte\\Apagado ordenado por USB/red};
\end{tikzpicture}
\caption{Topología física propuesta: disposición del rack de la oficina central.}
\label{fig:ce03-rack}
\end{figure}

\begin{table}[htbp]
\centering
\caption{Equipamiento físico propuesto (verificar compatibilidad con la guía de Broadcom/VMware antes de la compra).}
\label{tab:ce03-hardware}
\small
\begin{tabularx}{\textwidth}{L{2.4cm}L{4.4cm}X}
\toprule
\textbf{Equipo} & \textbf{Modelo propuesto} & \textbf{Especificación clave} \\
\midrule
ESXI-01 & Dell PowerEdge R250 (1U) & Xeon E-2378 (8 núcleos / 16 hilos), 64~GB DDR4 ECC (ampliable a 128~GB), 2 SSD SATA 1,92~TB en RAID~1 (PERC H355), 2 puertos 1~GbE + tarjeta Intel I350-T4, iDRAC. \\
FW01 & MikroTik RB5009UG+S+IN & 7 puertos 1~GbE, 1 de 2,5~GbE, 1 SFP+; RouterOS~7. \\
SW-CORE & MikroTik CRS326-24G-2S+RM & 24 puertos 1~GbE, 2 SFP+, VLAN 802.1Q, puerto espejo. \\
NAS-01 & Synology DS224+ & 2 discos de 4~TB en RAID~1 para respaldos locales. \\
UPS-01 & APC Smart-UPS 1500 VA rack 2U & Unos 1000~W; carga estimada de 240~W (unos 30 min de autonomía según la curva del fabricante). \\
Rack & Gabinete de 18U ventilado & Con PDU y organizadores. \\
\bottomrule
\end{tabularx}
\end{table}

\noindent El equipo actual (A8-7600 con 4~GB) no cumple los requisitos mínimos de ESXi~8
(8~GB de RAM y CPU compatible), por lo que se reutiliza como estación de pruebas.

\subsection{Diseño del hipervisor}

\begin{figure}[H]
\centering
\begin{tikzpicture}[x=1cm,y=1cm,scale=0.92,every node/.style={transform shape}]
  \node[cajaext,minimum width=1.4cm] (n0) at (-6,-5) {vmnic0};
  \node[cajaext,minimum width=1.4cm] (n1) at (-4.4,-5) {vmnic1};
  \node[cajaext,minimum width=1.4cm] (n2) at (0,-5) {vmnic2};
  \node[cajaext,minimum width=1.4cm] (n3) at (1.6,-5) {vmnic3};
  \node[cajaext,minimum width=1.4cm] (n4) at (6.0,-5) {vmnic4};
  \node[caja,minimum width=3.6cm,fill=upeublue!14] (vs0) at (-5.2,-3.4) {\textbf{vSwitch0} (gestión)\\activo/en espera};
  \node[caja,minimum width=3.6cm,fill=upeublue!14] (vs1) at (0.8,-3.4) {\textbf{vSwitch1} (datos)\\balanceo por puerto virtual};
  \node[caja,minimum width=2.4cm,fill=upeublue!14] (vs2) at (6.0,-3.4) {\textbf{vSwitch2}\\(espejo)};
  \draw[thick,upeublue] (n0) -- (vs0); \draw[thick,upeublue] (n1) -- (vs0);
  \draw[thick,upeublue] (n2) -- (vs1); \draw[thick,upeublue] (n3) -- (vs1); \draw[thick,upeublue] (n4) -- (vs2);
  \node[caja,font=\tiny] (pgm) at (-5.2,-1.9) {PG-GESTION (VLAN 10)\\VMkernel ESXi};
  \node[caja,font=\tiny] (pgs) at (-1.6,-1.9) {PG-SERVIDORES\\(VLAN 20)};
  \node[caja,font=\tiny] (pgd) at (0.8,-1.9) {PG-DMZ\\(VLAN 40)};
  \node[caja,font=\tiny] (pgsec) at (3.1,-1.9) {PG-SEGURIDAD\\(VLAN 50)};
  \node[cajasec,font=\tiny] (pgspan) at (6.0,-1.9) {PG-SPAN\\(VLAN 4095, promiscuo)};
  \draw[upeublue] (vs0) -- (pgm);
  \foreach \p in {pgs,pgd,pgsec} \draw[upeublue] (vs1.north) -- (\p.south);
  \draw[upeublue] (vs2) -- (pgspan);
  \node[caja,font=\tiny,text width=1.6cm] (app) at (-2.6,-0.2) {APP-01\\4 vCPU · 8 GB};
  \node[caja,font=\tiny,text width=1.6cm] (db) at (-0.7,-0.2) {DB-01\\4 vCPU · 8 GB};
  \node[caja,font=\tiny,text width=1.6cm] (px) at (1.2,-0.2) {PROXY-01\\2 vCPU · 2 GB};
  \node[caja,font=\tiny,text width=1.6cm] (hp) at (1.2,1.2) {HONEY-01\\1 vCPU · 1 GB};
  \node[caja,font=\tiny,text width=1.6cm] (mon) at (3.1,-0.2) {MON-01\\2 vCPU · 4 GB};
  \node[caja,font=\tiny,text width=1.6cm] (siem) at (3.1,1.2) {SIEM-01\\4 vCPU · 8 GB};
  \node[cajasec,font=\tiny,text width=1.6cm] (ids) at (6.0,-0.2) {IDS-01\\2 vCPU · 4 GB};
  \draw[upeublue] (pgs) -- (app); \draw[upeublue] (pgs) -- (db);
  \draw[upeublue] (pgd) -- (px); \draw[upeublue] (px) -- (hp);
  \draw[upeublue] (pgsec) -- (mon); \draw[upeublue] (mon) -- (siem);
  \draw[upeublue] (pgspan) -- (ids); \draw[upeublue,dashed] (ids.west) -- (mon.east);
  \begin{scope}[on background layer]
    \node[zona,fit=(vs0)(vs2)(app)(siem)(ids)(pgm),label={[zonatit]above:ESXI-01 — VMware ESXi 8}] {};
  \end{scope}
  \node[font=\tiny,text=upeugray] at (0.8,-5.7) {uplinks troncales 802.1Q hacia SW-CORE · vmnic4 al puerto espejo};
\end{tikzpicture}
\caption{Diseño interno del hipervisor: vSwitches, \emph{port groups}, NIC físicas y VM.}
\label{fig:ce03-esxi}
\end{figure}

\begin{table}[htbp]
\centering
\caption{Máquinas virtuales planeadas.}
\label{tab:ce03-vms}
\small
\begin{tabularx}{\textwidth}{L{1.9cm}C{0.9cm}C{1cm}L{2.1cm}L{2.4cm}X}
\toprule
\textbf{VM} & \textbf{vCPU} & \textbf{RAM} & \textbf{Disco} & \textbf{SO} & \textbf{Función} \\
\midrule
APP-01 & 4 & 8 GB & 60 + 150 GB & Ubuntu 24.04 + Docker & Backend FastAPI, panel, sitio, WhatsApp; disco de datos para evidencias \\
DB-01 & 4 & 8 GB & 120 GB & Ubuntu 24.04 & PostgreSQL 15, archivado WAL \\
PROXY-01 & 2 & 2 GB & 30 GB & Ubuntu 24.04 & Nginx + cloudflared (DMZ) \\
HONEY-01 & 1 & 1 GB & 20 GB & Ubuntu 24.04 & OpenCanary (señuelo DMZ) \\
MON-01 & 2 & 4 GB & 60 GB & Ubuntu 24.04 & Zabbix 7.0 LTS (SNMP a routers y OLT) \\
SIEM-01 & 4 & 8 GB & 150 GB & Ubuntu 24.04 & Wazuh (manager, indexer, dashboard) \\
IDS-01 & 2 & 4 GB & 40 GB & Ubuntu 24.04 & Suricata sobre el espejo del troncal \\
\midrule
\textbf{Total} & \textbf{19} & \textbf{35 GB} & \textbf{630 GB} & & \\
\bottomrule
\end{tabularx}
\end{table}

\begin{table}[htbp]
\centering
\caption{Red virtual del host.}
\label{tab:ce03-vswitch}
\small
\begin{tabularx}{\textwidth}{L{2.2cm}L{2.6cm}X}
\toprule
vSwitch0 & vmnic0 + vmnic1 & Gestión de ESXi (VMkernel, VLAN 10); NIC en activo/en espera para tolerar la falla de un puerto o cable. \\
vSwitch1 & vmnic2 + vmnic3 & Datos de las VM; troncal con VLAN 20, 40 y 50; balanceo por puerto virtual de origen. \\
vSwitch2 & vmnic4 & Recibe el puerto espejo del switch; \emph{port group} en VLAN 4095 con modo promiscuo solo para IDS-01. \\
Seguridad & --- & Modo promiscuo, cambios de MAC y transmisiones falsificadas en \emph{Rechazar} en todos los \emph{port groups} salvo PG-SPAN. \\
Licencia & --- & ESXi~8 en edición gratuita (disponible de nuevo desde 8.0 Update 3e) o licencia académica de Broadcom. La edición gratuita no expone la API de respaldo: se usan respaldos por script (ghettoVCB) y respaldos de aplicación. \\
\bottomrule
\end{tabularx}
\end{table}

% ======================================================================
\section{Arquitectura AWS --- diseño}

\subsection{Selección de región y justificación}

Se elige \textbf{América del Sur (São Paulo), \texttt{sa-east-1}}:
\begin{itemize}
  \item \textbf{Latencia:} es la región de AWS más cercana con servicios completos; la
  terminal SSH y el estado PPPoE en tiempo real del sistema pasan por el hub, así que
  la latencia se nota en la operación diaria. La zona local de Lima se evaluó, pero
  ofrece un catálogo reducido de servicios.
  \item \textbf{Normativa:} cualquier región implica flujo transfronterizo de datos
  personales; se documenta conforme a la Ley N.\textsuperscript{o}~29733 y su
  reglamento. En AWS solo residen una réplica y respaldos \textbf{cifrados}; la copia
  principal queda en la oficina.
  \item \textbf{Servicios:} EC2 Graviton (\texttt{t4g}), S3, Systems Manager,
  CloudWatch y VPC Flow Logs están disponibles en la región.
  \item \textbf{Costo:} es más cara que \texttt{us-east-1}; la diferencia es pequeña
  con dos instancias pequeñas y se acepta por la latencia. El costo se estimará con
  AWS Pricing Calculator antes de la implementación.
\end{itemize}

\subsection{Diseño de VPC}

\begin{table}[H]
\centering
\caption{Subredes de la VPC 10.50.0.0/16 (zona de disponibilidad \texttt{sa-east-1a}).}
\label{tab:ce03-subredes}
\small
\begin{tabularx}{\textwidth}{L{2.6cm}L{2.2cm}L{1.6cm}X}
\toprule
\textbf{Subred} & \textbf{CIDR} & \textbf{Tipo} & \textbf{Propósito} \\
\midrule
sn-publica & 10.50.1.0/24 & Pública & EC2-WGHUB con IP elástica (hub VPN, NAT para las privadas, honeypot) \\
sn-privada-app & 10.50.2.0/24 & Privada & EC2-DR (réplica y contingencia) \\
sn-privada-datos & 10.50.3.0/24 & Privada & Reservada para una base gestionada (RDS) futura \\
\bottomrule
\end{tabularx}
\end{table}

\begin{table}[htbp]
\centering
\caption{Tablas de rutas.}
\label{tab:ce03-rutas}
\small
\begin{tabularx}{\textwidth}{L{2.4cm}L{4.4cm}X}
\toprule
\textbf{Tabla} & \textbf{Destino} & \textbf{Destino de salto} \\
\midrule
rt-publica & 10.50.0.0/16 & local \\
 & 0.0.0.0/0 & Internet Gateway \\
rt-privada & 10.50.0.0/16 & local \\
 & 10.20.0.0/16, 10.10.0.0/24, 172.16.0.0/16 & ENI de EC2-WGHUB (vía VPN) \\
 & 0.0.0.0/0 & ENI de EC2-WGHUB (instancia NAT) \\
 & Lista de prefijos de S3 & Endpoint de puerta de enlace de S3 \\
\bottomrule
\end{tabularx}
\end{table}

\noindent Se usa la propia EC2-WGHUB como \textbf{instancia NAT} (comprobación de
origen/destino deshabilitada) en lugar de un NAT Gateway, cuyo cargo fijo mensual no
se justifica para el tráfico de actualizaciones de una sola instancia privada. Los
respaldos a S3 no pasan por el NAT: usan el endpoint de puerta de enlace, que no tiene
costo.

\begin{figure}[htbp]
\centering
\begin{tikzpicture}[x=1cm,y=1cm]
  \node[cajaext] (igw) at (-2.4,3.4) {Internet Gateway};
  \node[caja,text width=3.4cm] (hub) at (-2.4,1.2) {\textbf{EC2-WGHUB} t4g.medium\\10.50.1.10 + EIP\\WireGuard · NAT\\Suricata · Cowrie};
  \node[caja,text width=3.2cm] (dr) at (3.0,1.2) {\textbf{EC2-DR} t4g.medium\\10.50.2.10\\réplica PostgreSQL\\+ contenedores};
  \node[caja,text width=3.2cm,fill=upeugray!8] (dat) at (3.0,-1.3) {sn-privada-datos\\10.50.3.0/24 (reserva RDS)};
  \node[cajaext,text width=3.2cm] (s3) at (-2.4,-1.3) {Endpoint de S3\\$\rightarrow$ bucket de respaldos};
  \draw[flecha] (igw) -- node[font=\tiny,right]{rt-publica} (hub);
  \draw[flecha] (dr) -- node[font=\tiny,above,align=center]{rt-privada\\(VPN y NAT)} (hub);
  \draw[flecha,dashed] (dr.south west) -- (s3.north east);
  \begin{scope}[on background layer]
    \node[zona,fit=(hub),inner sep=6pt,label={[zonatit,font=\tiny\bfseries]below:sn-publica 10.50.1.0/24}] {};
    \node[zona,fit=(dr),inner sep=6pt,label={[zonatit,font=\tiny\bfseries]below:sn-privada-app 10.50.2.0/24}] {};
  \end{scope}
  \node[cajaext] (onp) at (-6.6,1.2) {Oficina\\y 13 zonas};
  \draw[tunel] (onp) -- node[font=\tiny,above]{WireGuard} (hub);
  \node[draw=upeublue,very thick,rounded corners,fit=(hub)(dr)(dat)(s3),inner sep=16pt,label={[zonatit]below:VPC 10.50.0.0/16 · sa-east-1}] {};
\end{tikzpicture}
\caption{Diseño de la arquitectura AWS (VPC).}
\label{fig:ce03-vpc}
\end{figure}

\subsection{Servicios EC2 planeados}

\begin{table}[H]
\centering
\caption{Instancias EC2 planeadas.}
\label{tab:ce03-ec2}
\small
\begin{tabularx}{\textwidth}{L{2.2cm}L{2.1cm}L{2.3cm}XL{1.8cm}}
\toprule
\textbf{Instancia} & \textbf{Tipo} & \textbf{Disco} & \textbf{Función} & \textbf{Subred} \\
\midrule
EC2-WGHUB & t4g.medium (2 vCPU, 4 GiB) & EBS gp3 20 GB cifrado & Hub WireGuard, nftables, instancia NAT, Suricata en línea y honeypot Cowrie & sn-publica \\
EC2-DR & t4g.medium (2 vCPU, 4 GiB) & EBS gp3 100 GB cifrado & Réplica en flujo de PostgreSQL y contenedores del sistema listos para activarse; se amplía a t4g.large al activar la contingencia & sn-privada-app \\
\bottomrule
\end{tabularx}
\end{table}

\subsection{Security Groups y NACL planeados}

\begin{table}[H]
\centering
\caption{Security Groups, NACL y roles de IAM.}
\label{tab:ce03-sg}
\small
\begin{tabularx}{\textwidth}{L{2.4cm}L{1.6cm}X}
\toprule
\textbf{Nombre} & \textbf{Alcance} & \textbf{Regla principal} \\
\midrule
sg-wghub & Instancia & Entrada: UDP 51820 desde 0.0.0.0/0 (zonas detrás de CGNAT); TCP 22 desde 0.0.0.0/0 \textbf{hacia Cowrie} (la administración real va por Systems Manager); todo desde sg-dr (NAT). \\
sg-dr & Instancia & Entrada: TCP 5432 solo desde 10.20.20.11/32 (DB-01 por VPN); TCP 8000 desde 10.20.10.0/24 (verificación en contingencia). Sin entrada desde internet. \\
NACL-publica & Subred & Permite UDP 51820, TCP 22 y puertos efímeros de retorno; deniega el resto. \\
NACL-privada & Subred & Permite tráfico desde 10.50.0.0/16, 10.20.0.0/16 y 10.10.0.0/24 y el retorno de puertos efímeros; deniega la entrada directa desde internet. \\
rol-ec2-hub & IAM & Systems Manager (sesión sin SSH) y agente de CloudWatch. \\
rol-ec2-dr & IAM & Systems Manager y escritura/lectura solo en el prefijo de respaldos del bucket. \\
svc-respaldo-oficina & IAM & Solo \texttt{PutObject} en el prefijo de respaldos (sin borrar). Cuenta raíz con MFA y sin uso diario; administradores con MFA. \\
\bottomrule
\end{tabularx}
\end{table}

\noindent El bucket de respaldos tiene bloqueo de acceso público, versionado, cifrado
del lado del servidor y una regla de ciclo de vida: 30 días en S3 Standard, luego S3
Glacier y eliminación al año. Las VPC Flow Logs se envían a CloudWatch para auditar
intentos de acceso (prueba PC-08 del diseño de red).

% ======================================================================
\section{Mecanismo de integración híbrida}

La integración es el túnel WireGuard diseñado en el \cref{chap:ce03-red}. Ambos
documentos usan las mismas redes y extremos:

\begin{table}[htbp]
\centering
\caption{Coherencia entre el diseño de red y el de centro de datos.}
\label{tab:ce03-coherencia}
\small
\begin{tabularx}{\textwidth}{L{3.6cm}L{4.6cm}X}
\toprule
\textbf{Elemento} & \textbf{Diseño de red} & \textbf{Diseño de centro de datos} \\
\midrule
Hub VPN & EC2-WGHUB, 10.10.0.1, \texttt{<EIP-HUB>}:51820 & EC2-WGHUB en sn-publica, 10.50.1.10 \\
Peer on-premise & FW01, 10.10.0.2; anuncia 10.20.\{10,20,50\}.0/24 & VM en VLAN 20 y 50 de ESXI-01 \\
Red anunciada por AWS & 10.50.0.0/16 & VPC 10.50.0.0/16 \\
Réplica de base de datos & Flujo DB-01 $\rightarrow$ EC2-DR, TCP 5432 & sg-dr permite 5432 solo desde 10.20.20.11 \\
Peers de zona & 10.10.0.(100+\textit{id}), anuncian 172.16.\textit{id}.0/28 & rt-privada enruta 172.16.0.0/16 al hub \\
\bottomrule
\end{tabularx}
\end{table}

% ======================================================================
\section{Matriz de componentes del centro de datos}

{\small
\begin{longtable}{L{2.6cm}L{2cm}L{2.2cm}L{4.4cm}L{1.6cm}}
\caption{Matriz consolidada on-premise + AWS (con recursos asignados).}\label{tab:ce03-componentes-dc}\\
\toprule
\textbf{Componente} & \textbf{Entorno} & \textbf{Tipo} & \textbf{Función y recursos} & \textbf{Criticidad} \\
\midrule
\endfirsthead
\toprule
\textbf{Componente} & \textbf{Entorno} & \textbf{Tipo} & \textbf{Función y recursos} & \textbf{Criticidad} \\
\midrule
\endhead
\bottomrule
\endlastfoot
ESXI-01 & On-premise & Hipervisor & Virtualización; 8C/16T, 64 GB, 1,92 TB RAID 1 & Alta \\
Datastore SSD & On-premise & Almacenamiento & Discos de las VM (RAID 1) & Alta \\
NAS-01 & On-premise & Almacenamiento & Respaldos locales; 4 TB RAID 1 & Media \\
FW01 & On-premise & Red & Firewall y peer VPN & Alta \\
SW-CORE & On-premise & Red & Troncal y espejo & Alta \\
UPS-01 & On-premise & Energía & 1500 VA, unos 30 min & Alta \\
APP-01 & On-premise & VM & Sistema; 4 vCPU, 8 GB, 210 GB & Alta \\
DB-01 & On-premise & VM & PostgreSQL; 4 vCPU, 8 GB, 120 GB & Alta \\
PROXY-01 & On-premise & VM & Publicación; 2 vCPU, 2 GB, 30 GB & Alta \\
HONEY-01 & On-premise & VM & Señuelo; 1 vCPU, 1 GB, 20 GB & Media \\
MON-01 & On-premise & VM & Monitoreo; 2 vCPU, 4 GB, 60 GB & Media \\
SIEM-01 & On-premise & VM & SIEM; 4 vCPU, 8 GB, 150 GB & Media \\
IDS-01 & On-premise & VM & IDS; 2 vCPU, 4 GB, 40 GB & Media \\
VPC principal & AWS & Red virtual & Aislamiento de recursos cloud & Alta \\
Internet Gateway & AWS & Red & Salida de sn-publica & Alta \\
Endpoint S3 & AWS & Red & Acceso privado a S3 & Media \\
EC2-WGHUB & AWS & Cómputo & Hub VPN; t4g.medium, 20 GB & Alta \\
EC2-DR & AWS & Cómputo & Contingencia; t4g.medium, 100 GB & Media \\
Bucket S3 de respaldos & AWS & Almacenamiento & Respaldos cifrados y versionados & Alta \\
CloudWatch / Flow Logs & AWS & Monitoreo & Métricas, alarmas y auditoría de red & Media \\
\end{longtable}
}

% ======================================================================
\section{Dimensionamiento y capacidad}

\begin{table}[H]
\centering
\caption{Recursos requeridos frente a la capacidad del host ESXI-01.}
\label{tab:ce03-capacidad}
\small
\begin{tabularx}{\textwidth}{L{2.4cm}L{2.6cm}L{2.6cm}X}
\toprule
\textbf{Recurso} & \textbf{Asignado (VM)} & \textbf{Capacidad} & \textbf{Uso y margen} \\
\midrule
CPU & 19 vCPU & 8 núcleos / 16 hilos & Relación 1,2:1 sobre hilos (2,4:1 sobre núcleos), por debajo del 3:1 habitual en cargas ligeras. \\
RAM & 35 GB + unos 4 GB del hipervisor & 64 GB & 61\,\% de uso, \textbf{39\,\% de margen}; ampliable a 128 GB. \\
Disco & 630 GB + unos 40 GB de \emph{swap} de VM & 1,92 TB útiles & 35\,\% de uso; queda espacio para instantáneas y para duplicar la base y las evidencias. \\
Red & Menos de 50 Mbps sostenidos & 4 $\times$ 1 GbE de datos/gestión & Holgado; la VPN de gestión usa menos de 2 Mbps. \\
\bottomrule
\end{tabularx}
\end{table}

\paragraph{Justificación de las cargas.} El sistema completo consume hoy entre 1,3 y
1,8~GB de RAM en reposo; los 8~GB de APP-01 cubren la compilación del frontend (1,5 a
2~GB puntuales), el techo de 512~MB del servicio de WhatsApp y el crecimiento. La base
de datos tiene 21 tablas con volúmenes pequeños (cientos de clientes, unas decenas de
miles de pagos e historial en 3 años, menos de 2~GB), por lo que 8~GB de RAM permiten
mantenerla en memoria; los 120~GB incluyen el archivado de WAL. Las evidencias crecen
unos 30~MB por cliente: unos 27~GB a 3 años, frente a 150~GB asignados. Wazuh y
Zabbix se dimensionan según los mínimos de sus fabricantes para menos de 50 agentes o
equipos.

\paragraph{AWS.} El hub solo transporta gestión (menos de 2~Mbps); \texttt{t4g.medium}
se elige por la RAM que necesitan Suricata con reglas ET Open (1 a 1,5~GB) y Cowrie, no
por la CPU. EC2-DR replica una base pequeña con 4~GiB y se amplía a \texttt{t4g.large}
solo cuando se activa la contingencia. Las imágenes del sistema (Python, Node,
PostgreSQL) son multiarquitectura, por lo que corren en Graviton (ARM).

\paragraph{Energía.} ESXI-01 (unos 150~W), FW01 y SW-CORE (unos 40~W) y NAS-01 (unos
25~W) suman unos 240~W con margen, dentro de la capacidad del UPS de 1500~VA.

% ======================================================================
\section{Consideraciones de disponibilidad y continuidad en el diseño}

\begin{itemize}
  \item \textbf{Redundancia local:} RAID~1 en el datastore y en el NAS, NIC en
  activo/en espera para la gestión y balanceo para los datos, UPS con apagado ordenado
  y encendido automático tras un corte (configuración de BIOS que ya se aplica hoy).
  El R250 tiene una sola fuente de poder: es un punto único de falla aceptado y
  compensado con la réplica en AWS.
  \item \textbf{Respaldos 3-2-1:} (1) datos en producción en ESXI-01; (2) copia local
  en NAS-01 (respaldos de VM con ghettoVCB y \texttt{pg\_dump} diario); (3) copia fuera
  de sitio en S3 (\texttt{pg\_dump} cifrado, archivado WAL cada 15 minutos y
  sincronización nocturna de evidencias).
  \item \textbf{Contingencia:} réplica en flujo de PostgreSQL hacia EC2-DR; ante la
  pérdida de la oficina se promueve la réplica, se levantan los contenedores en
  EC2-DR y se apunta el túnel de Cloudflare a esa instancia.
  \item \textbf{Configuración como código:} el \texttt{docker-compose.yml}, la plantilla
  de router y la configuración del hub se guardan en un repositorio privado (sin
  secretos) para reconstruir en minutos.
\end{itemize}

\begin{table}[htbp]
\centering
\caption{RTO/RPO objetivo preliminar por servicio (se validará en el semestre 2).}
\label{tab:ce03-rto}
\small
\begin{tabularx}{\textwidth}{L{2.8cm}C{1.6cm}C{1.3cm}C{1.2cm}X}
\toprule
\textbf{Servicio} & \textbf{Criticidad} & \textbf{RPO} & \textbf{RTO} & \textbf{Estrategia} \\
\midrule
Base de datos & Alta & 5 min & 1 h & Promover la réplica de EC2-DR; alternativamente, restaurar el respaldo con WAL (recuperación a un punto en el tiempo). \\
API, panel y app & Alta & 15 min & 2 h & Restaurar APP-01 desde NAS-01 o activar los contenedores en EC2-DR y redirigir el túnel. \\
Hub VPN & Alta & 24 h & 30 min & Relanzar desde la AMI y la plantilla de lanzamiento y reasignar la IP elástica; las zonas se reconectan solas. \\
Configuración de routers & Alta & 24 h & 1 h & Exportación nocturna (\texttt{/export}) al repositorio privado y restauración del \texttt{.rsc}. \\
Evidencias (fotos) & Media & 24 h & 4 h & Sincronización nocturna a S3 y NAS-01. \\
WhatsApp & Media & --- & 4 h & Recrear el contenedor y volver a vincular el número. \\
Monitoreo y SIEM & Media & 24 h & 8 h & Respaldo semanal de las VM. \\
Internet de clientes & Alta & --- & --- & No depende del centro de datos: cada zona sale por su ISP local. \\
\bottomrule
\end{tabularx}
\end{table}

% ======================================================================
\section{Conclusiones y hoja de ruta hacia la implementación}

El diseño pasa de un único equipo sin redundancia a un centro de datos híbrido con
cargas aisladas en VM, una copia de contingencia en AWS y respaldos 3-2-1, con un
costo acotado para una MYPE. Queda pendiente para el entregable
\emph{Implementación, gestión y control de centro de datos} del semestre~2:

\begin{enumerate}
  \item Adquirir e instalar ESXI-01, crear vSwitches, \emph{port groups} y las 7 VM.
  \item Migrar los contenedores actuales a APP-01, DB-01 y PROXY-01 sin pérdida de datos.
  \item Desplegar la VPC, EC2-WGHUB, EC2-DR, el bucket S3 y los roles de IAM con MFA.
  \item Configurar la réplica de PostgreSQL y los respaldos 3-2-1.
  \item Configurar el monitoreo (métricas del host ESXi y CloudWatch) con alertas de capacidad.
  \item Ejecutar y documentar al menos un ejercicio de recuperación completo, comparando
  el RTO/RPO obtenido con el objetivo de la \cref{tab:ce03-rto}.
\end{enumerate}

% ======================================================================
\section*{Referencias}
\addcontentsline{toc}{section}{Referencias}
{\small
\begin{enumerate}[label={[\arabic*]},leftmargin=2em]
  \item Broadcom / VMware, ``vSphere 8 Documentation: ESXi Installation and Setup, vSphere Networking''. [En línea]. Disponible: https://techdocs.broadcom.com
  \item Broadcom, ``Broadcom Compatibility Guide (servidores para ESXi)''. [En línea]. Disponible: https://compatibilityguide.broadcom.com
  \item Amazon Web Services, ``Amazon VPC User Guide'', ``Amazon EC2 User Guide'', ``IAM User Guide'' y ``Amazon S3 User Guide''. [En línea]. Disponible: https://docs.aws.amazon.com
  \item Amazon Web Services, ``AWS Well-Architected Framework'', 2024.
  \item ISO/IEC 27001:2022 e ISO/IEC 27002:2022 (controles 8.13 Respaldo de la información, 8.14 Redundancia, 5.23 Seguridad en servicios en la nube). Ginebra: ISO, 2022.
  \item PostgreSQL Global Development Group, ``PostgreSQL 15 Documentation: High Availability, Load Balancing, and Replication''.
  \item Congreso de la República del Perú, Ley N.\textsuperscript{o}~29733, Ley de Protección de Datos Personales.
\end{enumerate}
}
